% problem-000041 1 铝黄铜组成与超细铜制备
\section*{1. 铝黄铜组成与超细铜制备}
\textit{下列为分问。}

\subsection*{1-1}
1.000 g 铝⻩铜合⾦(设只含铜、锌、铝)与 0.100 mol L−1硫酸反应，在 $2 5 ^ { \circ } \mathrm { C }$ 和 101.325 kPa 下测得放出的⽓体的体积为149.3 mL。将相同质量的该合⾦溶于⾜量热浓硫酸，在相同温度和压强下测得放出的⽓体的体积为 $4 1 1 . 1 \mathrm { m L } _ { \circ }$ 计算此铝⻩铜合⾦中各组分的质量分数。

\subsection*{1-2}
向 $\mathrm { [ C u ( N H _ { 3 } ) _ { 4 } ] S O _ { 4 } }$ ⽔溶液中通⼊ $\mathrm { S O } _ { 2 }$ ⽓体⾄溶液呈微酸性，析出⽩⾊沉淀 $\mathrm { C u N H _ { 4 } S O _ { 3 \circ } }$ $\mathrm { C u N H _ { 4 } S O _ { 3 } }$ 与⾜量的硫酸混合并微热，得到⾦属 $\mathrm { C u }$ 等物质，本法制得的Cu呈超细粉末状，有重要⽤途。

\subsection*{1-2}
-1 写出⽣成 $\mathrm { C u N H _ { 4 } S O _ { 3 } }$ 的反应⽅程式。

\subsection*{1-2}
-2 写出 $\mathrm { C u N H _ { 4 } S O _ { 3 } }$ 与 H $\mathrm { _ { 2 S O _ { 4 } } }$ 作⽤的反应⽅程式，若反应在敞开反应器中进⾏，计算反应物中的Cu元素变成超细粉末Cu的转化率。

\subsection*{1-2}
-3 若反应在密闭容器中进⾏，且酸量充⾜，计算反应物中的Cu元素变成超细粉末Cu的转化率。并对此做出解释。
